\section{Related Work}
\label{sec:related-work}

\subsection{Multi-agent critique and model revision}

Multi-agent LLM methods use interaction among model instances to improve a
final response. Multi-agent debate, for example, allows agents to exchange
arguments and revise their answers over several rounds, producing gains on
some factual and reasoning tasks \cite{du2024debate}. A Solver--Critic system
is a simpler instance of the same idea: one model proposes an answer and a
second model evaluates it before revision. These approaches are motivated by
the possibility that one agent can expose errors missed by another.

Revision is not automatically equivalent to correction. Huang et
al.\ show that LLMs can struggle to improve reasoning through intrinsic
self-correction without reliable external feedback
\cite{huang2024selfcorrect}. Conversely, models may accept external signals
for reasons unrelated to their correctness. Work on sycophancy documents cases
in which models adapt responses to a user's stated view even when doing so
reduces truthfulness \cite{sharma2024sycophancy}. Together, these findings
motivate separating three questions: whether feedback is correct, whether the
Solver adopts it, and whether adoption improves the final answer. Our primary
behavioral target concerns adoption rather than correctness, while correctness
is retained for stratified analysis.

\subsection{Sparse autoencoders for language-model representations}

Transformer residual streams are high-dimensional and superposed: a single
dimension need not correspond to a single human-interpretable concept. SAEs
learn an overcomplete dictionary that reconstructs dense activations from a
sparse set of latent features. Cunningham et al.\ demonstrated that SAE
features can be substantially more interpretable than raw neuron activations
\cite{cunningham2023sparse}. Subsequent work studied how SAE design and scale
affect sparsity, reconstruction quality, and downstream interpretability
\cite{gao2025scaling}.

SAE features remain learned model components rather than guaranteed semantic
units. Reconstruction error, inactive features, correlated latents, and
feature splitting can all complicate interpretation. For that reason, this
study compares multiple SAE variants and measures both reconstruction behavior
and downstream predictive utility. A general pretrained SAE tests whether
broadly learned features transfer to feedback processing, whereas
interaction-trained SAEs test whether training on the experiment's activation
distribution better captures the relevant structure.

\subsection{Probing and causal validation}

Linear probes test whether information can be decoded from an intermediate
representation \cite{alain2016probes}. Their simplicity makes coefficients
useful for ranking candidate features, but probe performance alone does not
show that the model uses those features to produce its behavior. Probes may
exploit correlated information or properties of the dataset rather than a
causal mechanism. Control tasks, held-out evaluation, and careful capacity
selection are therefore necessary when interpreting probe results
\cite{hewitt2019control}.

Feature intervention provides a stronger test. Work on sparse feature
circuits demonstrates how sparse features can support both circuit discovery
and targeted model editing \cite{marks2025circuits}. Following this general
principle, our analysis uses probe coefficients only to nominate candidate SAE
features. We then replay held-out Solver interactions while ablating selected
features and compare the resulting adoption rate with no-intervention and
random-feature controls. This progression from decoding to intervention is
intended to distinguish features that merely describe feedback uptake from
features that help produce it.

\subsection{Multi-hop question answering as an experimental setting}

MuSiQue constructs multi-hop questions by composing connected single-hop
questions and filtering shortcuts that allow models to bypass the intended
reasoning chain \cite{trivedi2022musique}. It is useful for the present study
because it supplies question identifiers, supporting passages, known answers,
and answer aliases at sufficient scale for repeated Solver--Critic
interactions. Its range of two- to four-hop questions also creates variation
in difficulty. At the same time, exact or normalized string matching cannot
resolve every semantic equivalence. Our labeling procedure therefore uses
question-local aliases and conservative text rules, and excludes cases whose
adoption status cannot be determined reliably.
